\newpage
\section{Introduction}
\label{sec:introduction}

The general topic of our project is conservation of linear momentum in
collisions. When two objects hit each other, the forces between them are
internal forces of the two-object system. During the very short contact
they are also much bigger than any outside force, so the total momentum
stays the same before and after the collision:
\begin{equation}
  m\vect v_{d0}+M\vect V_0 = m\vect v_{d1}+M\vect V_1 ,
  \label{eq:intro-momentum}
\end{equation}
where $m$ and $M$ are the masses of the debris fragment and the satellite,
$\vect v_d$ and $\vect V$ are their velocities, and the subscripts 0 and 1
mean before and after the hit. Equation~\eqref{eq:intro-momentum} is true
for every kind of collision. What changes from one kind to another is the
kinetic energy: it is kept in an elastic collision and partly lost in an
inelastic one \citep{halliday2014, ling2016}.

Our specific title is \textit{Space Debris Impact on Solar Panels}. Low
Earth orbit is full of small fragments left over from old explosions and
collisions, and every new collision makes more of them. This chain
reaction was first described by \citet{kessler1978}. Solar panels are
usually the biggest surfaces on a satellite, so they get hit the most. A
real example happened on 23 August 2016, when a particle about 1~cm in
size hit a solar panel of the Sentinel-1A satellite. The satellite lost
part of its power, and its orbit and orientation changed suddenly
\citep{krag2017}. This event shows the two effects we study: a change in
the satellite's velocity and a change in its spin.

Both effects can be explained with basic physics. The fragment gives part
of its momentum to the satellite (Eq.~\eqref{eq:intro-momentum}). Because
the panel is far from the satellite's centre of mass, the same push also
creates a torque, so the satellite starts to rotate. This is described by
conservation of angular momentum. Kinetic energy is different, because it
grows with $v^2$. A fragment moving at $10\unit{km\,s^{-1}}$ carries a lot
of energy, and in a completely inelastic hit most of it turns into heat
and damage on the panel.

To solve the problem, we model the satellite as a rigid body and the
fragment as a point mass. We then solve the two limiting collisions from
the course exactly. In the elastic case kinetic energy is kept, and in the
completely inelastic case the fragment sticks to the panel and both move
as one rigid body. Our Python program calculates the satellite's velocity
change, its spin and the energy lost for three scenarios and several
parameter sweeps, and compares the velocity change with the orbital speed.
We check the code in two ways. First, the momentum, angular momentum and
energy before and after the hit must be equal up to round-off error.
Second, a head-on hit through the centre of mass must give the textbook
1-D results: with equal masses, the target should move at $v_0$ in the
elastic case and at $v_0/2$ in the completely inelastic case
\citep{ling2016}. At the end we compare our results with what was
measured for Sentinel-1A \citep{krag2017}.
